\section{Refrescos S.A.}


\subsection{Presentación del caso}

La compañía \textbf{REFRESCOS S.A.}, dedicada a la fabricación y distribución de bebidas refrescantes, ha detectado una modificación en la demanda de sus clientes y plantea estudiar el interés de cerrar alguna de sus unidades productivas.  
Se produce y distribuye un único producto.  
La empresa cuenta con tres líneas de fabricación en la actualidad: \textbf{Madrid}, \textbf{Toledo} y \textbf{Segovia}.  
Dispone de cinco almacenes de distribución: Madrid (mismo emplazamiento que la línea), Segovia (mismo emplazamiento que la línea) y tres en Toledo (Toledo 1, Toledo 2 y Toledo 3).  
En cada línea se pueden realizar dos turnos diarios.  

Los costes fijos de funcionamiento de las líneas son:  
Madrid: 90.000 €; Toledo: 90.000 €; Segovia: 12.000 €.  

Los costes fijos por turno en cada línea son:

\begin{center}
\begin{tabular}{lcc}
\toprule
Línea & Turno 1 (€) & Turno 2 (€) \\
\midrule
Madrid & 24.000 & 30.000 \\
Toledo & 12.000 & 15.000 \\
Segovia & 6.000 & 7.500 \\
\bottomrule
\end{tabular}
\end{center}

Los costes variables de transporte (en €/caja) son:

\begin{center}
\begin{tabular}{lcl}
Madrid $\rightarrow$ Toledo 1 & : & 0.12 \\
Madrid $\rightarrow$ Toledo 2 & : & 0.15 \\
Madrid $\rightarrow$ Toledo 3 & : & 0.18 \\
Madrid $\rightarrow$ Segovia & : & 0.09 \\
Toledo $\rightarrow$ Toledo 2 & : & 0.06 \\
Toledo $\rightarrow$ Toledo 3 & : & 0.09 \\
Toledo $\rightarrow$ Segovia & : & 0.21 \\
Segovia $\rightarrow$ Toledo 2 & : & 0.24 \\
Segovia $\rightarrow$ Toledo 3 & : & 0.27 \\
\end{tabular}
\end{center}

Las líneas tienen la siguiente capacidad de producción:

\begin{itemize}
    \item Madrid: 150 horas/mes por turno. Produce 1000 cajas en 1,2 horas.
    \item Toledo: 140 horas/mes por turno. Produce 1000 cajas en 2,8 horas.
    \item Segovia: 140 horas/mes por turno. Produce 1000 cajas en 2 horas.
\end{itemize}

La demanda mensual es la siguiente:

\begin{center}
\begin{tabular}{lc}
\toprule
Almacén & Demanda (miles de cajas/mes) \\
\midrule
Madrid & 200 \\
Toledo 1 & 30 \\
Toledo 2 & 10 \\
Toledo 3 & 10 \\
Segovia & 30 \\
\bottomrule
\end{tabular}
\end{center}




\subsection{Formulación explícita}

\paragraph{Variables}\mbox{}\\

\begin{tabular}{lcl}
$x_{la}$ & : & miles de cajas transportadas desde la línea $l$ al almacén $a$ \\
$\alpha_l$ & : & 1 si la línea $l$ está abierta; 0 en caso contrario \\
$\beta_{lt}$ & : & 1 si el turno $t$ de la línea $l$ está abierto; 0 en caso contrario \\
\end{tabular}

\paragraph{Función objetivo}\mbox{}\\

Minimizar los costes totales:

\begin{align}
\mbox{min. } z = & \; 0.12x_{12} + 0.15x_{13} + 0.18x_{14} + 0.09x_{15} \notag\\
& + 0.21x_{21} + 0.06x_{23} + 0.09x_{24} + 0.21x_{25} \notag\\
& + 0.09x_{31} + 0.21x_{32} + 0.24x_{33} + 0.27x_{34} \notag\\
& + 90000\alpha_1 + 90000\alpha_2 + 12000\alpha_3 \notag\\
& + 24000\beta_{11} + 30000\beta_{12} + 12000\beta_{21} + 15000\beta_{22} + 6000\beta_{31} + 7500\beta_{32}
\end{align}

\paragraph{Restricciones}\mbox{}\\

Producción máxima disponible en cada línea:  
(La capacidad se calcula como $H_{tl} \times P_l$ para cada turno abierto)

\begin{align}
x_{11}+x_{12}+x_{13}+x_{14}+x_{15} &\le (150\times0.8333)\beta_{11} + (150\times0.8333)\beta_{12} = 125\,\beta_{11} + 125\,\beta_{12} \\
x_{21}+x_{22}+x_{23}+x_{24}+x_{25} &\le (140\times0.3571)\beta_{21} + (140\times0.3571)\beta_{22} = 50\,\beta_{21} + 50\,\beta_{22} \\
x_{31}+x_{32}+x_{33}+x_{34}+x_{35} &\le (140\times0.5)\beta_{31} + (140\times0.5)\beta_{32} = 70\,\beta_{31} + 70\,\beta_{32}
\end{align}

Demanda en los almacenes:

\begin{align}
x_{11}+x_{21}+x_{31} &\ge 200 \\
x_{12}+x_{22}+x_{32} &\ge 30 \\
x_{13}+x_{23}+x_{33} &\ge 10 \\
x_{14}+x_{24}+x_{34} &\ge 10 \\
x_{15}+x_{25}+x_{35} &\ge 30
\end{align}

Relación entre líneas y turnos:

\begin{align}
\beta_{11} &= \alpha_1 \\
\beta_{21} &= \alpha_2 \\
\beta_{31} &= \alpha_3 \\
\beta_{11} &\ge \beta_{12} \\
\beta_{21} &\ge \beta_{22} \\
\beta_{31} &\ge \beta_{32}
\end{align}

Restricciones de no negatividad e integralidad:

\begin{align}
x_{la} &\ge 0 \\
\alpha_l,\,\beta_{lt} &\in \{0,1\}
\end{align}

\paragraph{Nota de unidades}\mbox{}\\
Todas las variables de flujo $x_{la}$ y demandas $D_a$ están en \textbf{miles de cajas}.  
Los costes de transporte $CT_{la}$ están en \textbf{euros por mil cajas}.  
En la restricción de capacidad, $H_{tl}$ (horas por turno en la línea $l$) y $P_l$ (miles de cajas por hora en la línea $l$) generan capacidades en miles de cajas coherentes con $x_{la}$.
**

\subsection{Formulación generalizada}

\paragraph{Conjuntos}\mbox{}\\
\begin{tabular}{lcl}
$\mathcal{L}$ & : & líneas de producción\\
$\mathcal{A}$ & : & almacenes de distribución\\
$\mathcal{T}$ & : & turnos por línea\\
\end{tabular}

\paragraph{Parámetros}\mbox{}\\
\begin{tabular}{lcl}
$CT_{la}$ & : & coste de transporte por mil cajas desde $l$ a $a$\\
$CF_l$ & : & coste fijo por abrir la línea $l$\\
$CT_{lt}$ & : & coste fijo por abrir el turno $t$ en la línea $l$\\
$H_{tl}$ & : & horas disponibles en el turno $t$ de la línea $l$\\
$P_l$ & : & productividad (miles de cajas/hora) de la línea $l$\\
$D_a$ & : & demanda (miles de cajas) del almacén $a$\\
\end{tabular}

\paragraph{Variables}\mbox{}\\
\begin{tabular}{lcl}
$x_{la}$ & : & miles de cajas enviadas de $l$ a $a$\\
$\alpha_l$ & : & 1 si la línea $l$ está abierta\\
$\beta_{lt}$ & : & 1 si el turno $t$ de la línea $l$ está abierto\\
\end{tabular}

\paragraph{Modelo completo}\mbox{}\\
\begin{align}
\mbox{min. } z \;=\;& \sum_{l\in\mathcal{L}}\sum_{a\in\mathcal{A}} CT_{la}\,x_{la} \;+\; \sum_{l\in\mathcal{L}} CF_l\,\alpha_l \;+\; \sum_{l\in\mathcal{L}}\sum_{t\in\mathcal{T}} CT_{lt}\,\beta_{lt} \notag\\[1mm]
\mbox{s. a.:}\quad
& \sum_{a\in\mathcal{A}} x_{la} \;\le\; \sum_{t\in\mathcal{T}} H_{tl}\,P_l\,\beta_{lt} && \forall l\in\mathcal{L} \notag\\
& \sum_{l\in\mathcal{L}} x_{la} \;\ge\; D_a && \forall a\in\mathcal{A} \notag\\
& \beta_{l1} = \alpha_l && \forall l\in\mathcal{L} \notag\\
& \beta_{l2} \le \beta_{l1} && \forall l\in\mathcal{L} \notag\\
& x_{la}\ge 0 && \forall l\in\mathcal{L},\,a\in\mathcal{A} \notag\\
& \alpha_l,\,\beta_{lt}\in\{0,1\} && \forall l\in\mathcal{L},\,t\in\mathcal{T} \notag
\end{align}

Los valores de los parámetros, de acuerdo con la información proporcionada son os siguientes.

\begin{table}[!h]
\centering
\caption{parámetros del caso \textit{Refrescos S.A.}}
\begin{tabular}{lcl}
\toprule
\textbf{Parámetro} & \textbf{Símbolo} & \textbf{Valor} \\
\midrule
\textbf{Demanda (miles de cajas/mes)} & & \\
\quad Madrid & $D_{\text{Madrid}}$ & 200 \\
\quad Toledo 1 & $D_{\text{Toledo1}}$ & 30 \\
\quad Toledo 2 & $D_{\text{Toledo2}}$ & 10 \\
\quad Toledo 3 & $D_{\text{Toledo3}}$ & 10 \\
\quad Segovia & $D_{\text{Segovia}}$ & 30 \\
\midrule
\textbf{Costes fijos de apertura de línea (€/mes)} & $CF_l$ & \\
\quad Madrid & $CF_{\text{Madrid}}$ & 90{,}000 \\
\quad Toledo & $CF_{\text{Toledo}}$ & 90{,}000 \\
\quad Segovia & $CF_{\text{Segovia}}$ & 12{,}000 \\
\midrule
\textbf{Costes fijos por turno (€/mes)} & $CT_{lt}$ & \\
\quad Madrid, turno 1 & $CT_{\text{Madrid},1}$ & 24{,}000 \\
\quad Madrid, turno 2 & $CT_{\text{Madrid},2}$ & 30{,}000 \\
\quad Toledo, turno 1 & $CT_{\text{Toledo},1}$ & 12{,}000 \\
\quad Toledo, turno 2 & $CT_{\text{Toledo},2}$ & 15{,}000 \\
\quad Segovia, turno 1 & $CT_{\text{Segovia},1}$ & 6{,}000 \\
\quad Segovia, turno 2 & $CT_{\text{Segovia},2}$ & 7{,}500 \\
\midrule
\textbf{Horas disponibles por turno} & $H_{tl}$ & \\
\quad Madrid & $H_{\text{Madrid}}$ & 150 \\
\quad Toledo & $H_{\text{Toledo}}$ & 140 \\
\quad Segovia & $H_{\text{Segovia}}$ & 140 \\
\midrule
\textbf{Productividad (miles de cajas/hora)} & $P_l$ & \\
\quad Madrid & $P_{\text{Madrid}}$ & 0.8333 \\
\quad Toledo & $P_{\text{Toledo}}$ & 0.3571 \\
\quad Segovia & $P_{\text{Segovia}}$ & 0.5000 \\
\midrule
\textbf{Costes de transporte (€/mil cajas)} & $CT_{la}$ & \\
\quad Madrid $\rightarrow$ Toledo 1 & $CT_{\text{Madrid,Toledo1}}$ & 120 \\
\quad Madrid $\rightarrow$ Toledo 2 & $CT_{\text{Madrid,Toledo2}}$ & 150 \\
\quad Madrid $\rightarrow$ Toledo 3 & $CT_{\text{Madrid,Toledo3}}$ & 180 \\
\quad Madrid $\rightarrow$ Segovia & $CT_{\text{Madrid,Segovia}}$ & 90 \\
\quad Toledo $\rightarrow$ Toledo 2 & $CT_{\text{Toledo,Toledo2}}$ & 60 \\
\quad Toledo $\rightarrow$ Toledo 3 & $CT_{\text{Toledo,Toledo3}}$ & 90 \\
\quad Toledo $\rightarrow$ Segovia & $CT_{\text{Toledo,Segovia}}$ & 210 \\
\quad Segovia $\rightarrow$ Toledo 2 & $CT_{\text{Segovia,Toledo2}}$ & 240 \\
\quad Segovia $\rightarrow$ Toledo 3 & $CT_{\text{Segovia,Toledo3}}$ & 270 \\
\bottomrule
\end{tabular}
\end{table}




% \subsection{Presentación del caso}

% La compañía \textbf{REFRESCOS S.A.}, dedicada a la fabricación y distribución de bebidas refrescantes, ha detectado una modificación en la demanda de sus clientes y plantea estudiar el interés de cerrar alguna de sus unidades productivas.  
% Se produce y distribuye un único producto.  
% La empresa cuenta con tres líneas de fabricación en la actualidad: \textbf{Madrid}, \textbf{Toledo} y \textbf{Segovia}.  
% Dispone de cinco almacenes de distribución: Madrid (mismo emplazamiento que la línea), Segovia (mismo emplazamiento que la línea) y tres en Toledo (Toledo 1, Toledo 2 y Toledo 3).  
% En cada línea se pueden realizar dos turnos diarios.  

% Los costes fijos de funcionamiento de las líneas son:  
% Madrid: 90.000 €; Toledo: 90.000 €; Segovia: 12.000 €.  

% Los costes fijos por turno en cada línea son:

% \begin{center}
% \begin{tabular}{lcc}
% \toprule
% Línea & Turno 1 (€) & Turno 2 (€) \\
% \midrule
% Madrid & 24.000 & 30.000 \\
% Toledo & 12.000 & 15.000 \\
% Segovia & 6.000 & 7.500 \\
% \bottomrule
% \end{tabular}
% \end{center}

% Los costes variables de transporte (en €/caja) son:

% \begin{center}
% \begin{tabular}{lcl}
% Madrid $\rightarrow$ Toledo 1 & : & 0.12 \\
% Madrid $\rightarrow$ Toledo 2 & : & 0.15 \\
% Madrid $\rightarrow$ Toledo 3 & : & 0.18 \\
% Madrid $\rightarrow$ Segovia & : & 0.09 \\
% Toledo $\rightarrow$ Toledo 2 & : & 0.06 \\
% Toledo $\rightarrow$ Toledo 3 & : & 0.09 \\
% Toledo $\rightarrow$ Segovia & : & 0.21 \\
% Segovia $\rightarrow$ Toledo 2 & : & 0.24 \\
% Segovia $\rightarrow$ Toledo 3 & : & 0.27 \\
% \end{tabular}
% \end{center}

% Las líneas tienen la siguiente capacidad de producción:

% \begin{itemize}
%     \item Madrid: 150 horas/mes por turno. Produce 1000 cajas en 1,2 horas.
%     \item Toledo: 140 horas/mes por turno. Produce 1000 cajas en 2,8 horas.
%     \item Segovia: 140 horas/mes por turno. Produce 1000 cajas en 2 horas.
% \end{itemize}

% La demanda mensual es la siguiente:

% \begin{center}
% \begin{tabular}{lc}
% \toprule
% Almacén & Demanda (miles de cajas/mes) \\
% \midrule
% Madrid & 200 \\
% Toledo 1 & 30 \\
% Toledo 2 & 10 \\
% Toledo 3 & 10 \\
% Segovia & 30 \\
% \bottomrule
% \end{tabular}
% \end{center}


% \subsection{Formulación generalizada}

% \paragraph{Conjuntos}\mbox{}\\

% \begin{tabular}{lcl}
% $\mathcal{L}$ & : & líneas de producción \\
% $\mathcal{A}$ & : & almacenes de distribución \\
% $\mathcal{T}$ & : & turnos disponibles por línea \\
% \end{tabular}

% \paragraph{Parámetros}\mbox{}\\

% \begin{tabular}{lcl}
% $C_{la}$ & : & coste de transporte por caja desde la línea $l$ al almacén $a$ \\
% $CF_l$ & : & coste fijo por abrir la línea $l$ \\
% $CT_{lt}$ & : & coste fijo por abrir el turno $t$ de la línea $l$ \\
% $P_{lt}$ & : & capacidad de producción (cajas) de la línea $l$ en el turno $t$ \\
% $D_a$ & : & demanda (cajas) del almacén $a$ \\
% \end{tabular}

% \paragraph{Variables}\mbox{}\\

% \begin{tabular}{lcl}
% $x_{la}$ & : & número de cajas transportadas desde la línea $l$ al almacén $a$ \\
% $\alpha_l$ & : & 1 si la línea $l$ está abierta \\
% $f_{lt}$ & : & 1 si el turno $t$ de la línea $l$ está abierto \\
% \end{tabular}

% \paragraph{Función objetivo}\mbox{}\\

% \begin{equation}
% \mbox{min. } z = \sum_{l \in \mathcal{L}}\sum_{a \in \mathcal{A}} C_{la}x_{la} + \sum_{l \in \mathcal{L}} CF_l \alpha_l + \sum_{l \in \mathcal{L}}\sum_{t \in \mathcal{T}} CT_{lt} f_{lt}
% \end{equation}

% \paragraph{Restricciones}\mbox{}\\

% Capacidad de producción:

% \begin{equation}
% \sum_{a \in \mathcal{A}} x_{la} \le \sum_{t \in \mathcal{T}} P_{lt} f_{lt} \quad \forall l \in \mathcal{L}
% \end{equation}

% Satisfacción de la demanda:

% \begin{equation}
% \sum_{l \in \mathcal{L}} x_{la} \ge D_a \quad \forall a \in \mathcal{A}
% \end{equation}

% Relación entre líneas y turnos:

% \begin{align}
% f_{l1} &= \alpha_l \quad \forall l \in \mathcal{L} \\
% f_{l1} &\ge f_{l2} \quad \forall l \in \mathcal{L}
% \end{align}

% No negatividad e integralidad:

% \begin{align}
% x_{la} &\ge 0 \quad \forall l \in \mathcal{L}, a \in \mathcal{A} \\
% \alpha_l, f_{lt} &\in \{0,1\}
% \end{align}
